This thesis introduced LIRA (Latent Invariant Regime Analysis), a framework for
finding hidden causal regimes in black hole accretion time series. The method uses
the power-law physics of the Fundamental Plane and models the observations with a
log-linear structural equation. Black hole mass is included in the intercept, while
four spectral-state variables describe changes over time. LIRA finds regime
boundaries through recursive partitioning, a lagged Chow $F$-test, and the Bayesian
Information Criterion. DYNOTEARS then estimates the causal graph inside each
regime. The results are checked with refutation tests, sensitivity analysis, and graph
falsification.

The synthetic experiments produced strong results. LIRA found the true boundary
within four observations out of ten thousand. It also found the correct number of
regimes. It recovered the structural coefficients to within $0.017$ of their true
values.

The isolated variable in each regime was identified correctly. Clustering methods
that ignored causal structure performed poorly. A pooled model also gave an averaged
structure that matched neither regime.

The two-gate design was useful. The Chow test alone would have created extra splits.
The BIC gate rejected those unnecessary splits.

Several limitations remain. The most important one is that every result in this
thesis was obtained from synthetic data. The synthetic experiments show that LIRA
works when the data follow the assumptions it was designed for. They do not show that
it works on real black hole observations.

Real X-ray light curves contain gaps. They also contain instrumental artefacts and
sources of variability that the generator does not reproduce. The pipeline may behave
differently once it is applied to them.

The simulation also uses the same linear and piecewise-stationary structure that LIRA
assumes during recovery. The tests are therefore favourable to the method by
construction.

Beyond the data itself, the model assumes linear and piecewise-stationary
relationships. Real accretion systems may change in nonlinear and continuous ways
instead. The identifiability result also depends on the collider structure used in
this thesis. It may not hold for a model with causal chains between the spectral
variables.

The sensitivity analysis showed that the weaker coupling in the Soft State supports
weaker causal claims than the Hard State.

The next step is to move from synthetic validation to real data. This is the single
most important piece of future work. Real observations are the only way to test
whether the regimes LIRA finds correspond to physical accretion states. They also
show whether the regimes are artefacts of the simulation.

The recommended sequence mirrors the phase structure used in this thesis. LIRA should
first be applied to well-monitored sources with dense coverage. GRS 1915+105 and
Cygnus X-1 are good candidates. Their known state transitions provide an external
check on the discovered boundaries.

The analysis should then be extended to the WATCHDOG catalogue. It offers a broader
sample of sources and longer baselines.

In both cases, the discovered regimes should be compared against independently
identified Hard and Soft States. The recovered causal graphs should be checked
against the physical expectations discussed in Chapter 2.

Any disagreement between the synthetic and real-data results should be treated as a
finding in its own right. It would locate exactly where the simulation assumptions
break down.

Other useful steps include testing nonlinear mechanisms. LIRA can also be compared
with CASTOR on the same benchmark. Future work can also allow causal structure to
change gradually rather than at fixed boundaries.

Together, these extensions would address the gap between the controlled synthetic
validation reported here and the noisier, less predictable behaviour of real
accretion systems.
